% Notes of the Montreal seminar, 1970 (cote 129, typed part): archive cote 129, site page 32 (PDF page 33).
% typescript. Draft compared: none.
% AI-transcribed by Claude Opus 5.5 (claude-opus-5-5, adaptive thinking, effort high), Message Batches API, batch mtl129_full (msgbatch_01AGnDKHwy6ATCM7ehUa6263), 2026-11-13; not yet checked by eye.
\scanpage{129/32}
\mspage{\hand{II.4}}

Si l'on suppose que  S  est un schéma sur  $\mathbb{F}_p$ , on a défini
(I.5) le schéma en groupes  $G^{(p/S)}$  et les morphismes de Frobenius rela-%
tif  $\mathbb{F}_{G/S}$  et de Verschiebung  $\mathbb{V}_G$ . Si  G  est fini localement libre sur
S , $G^{(p/S)}$  l'est aussi et on montre alors (SGA 3, VII$_A$, 4.3.3) que la
dualité de Cartier \uline{échange les morphismes de Frobenius relatif et de}
\uline{Verschiebung}, i.e. on a :
\[ (\mathbb{V}_G)^* = \mathbb{F}_{G^*/S} \quad ; \quad (\mathbb{F}_{G/S})^* = \mathbb{V}_{G^*} \quad . \]

1.5.   \uline{Exemples}

a) Soit  $\underline{G}_S$  le groupe fini sur  S  attaché au groupe fini
(abstrait)  G  (1.3.1). On vérifie aisément que l'algèbre affine de  $(\underline{G}_S)^*$
est isomorphe à  $\underline{O}_S[G]$ , la $\underline{O}_S$-algèbre du groupe  G , l'application dia-%
gonale provenant de l'application diagonale  $G \rightarrow G \times G$ . Si  $G = \mathbb{Z}/n\mathbb{Z}$ ,
on trouve
\[ (\mathbb{Z}/n\mathbb{Z})^*_S \simeq \mu_{n|S} \quad , \]
\margin{$\times$} qui résulte aussi \deleted{diversement} \inserted{directement} des définitions.

b) En caractéristique  p , on montre que l'on a
\[ (\alpha_{p|S})^* \simeq \alpha_{p|S} \quad . \]
\margin{$\times$ \quad $\alpha_{p|S}$}

1.6.   \uline{Quotients de groupes finis localement libres}

Soit  $\mu : G' \rightarrow G$  un monomorphisme de groupes finis localement
libres sur  S , et soit  $G'' = G/G'$  le faisceau quotient (pour la topolo-%
gie f.p.p.f.). On a donc une suite exacte de faisceaux
\margin{$\pi$}
\[ 0 \rightarrow G' \xrightarrow{\mu} G \xrightarrow{\deleted{p}\inserted{\pi}} G'' \rightarrow 0 \quad . \]

\margin{32}
